\hypertarget{differntial-algebraic-equations}{%
\subsection{Differential-Algebraic
Equations}\label{sec:daes}}

% \newglossaryentry{fce}
{
    type=matprop,
    name=$f'_{ce}$,
    description={unconfined concrete compressive strength}
}
\newglossaryentry{ec0}
{
    type=matprop,
    name=$\epsilon_{c0}$,
    description={unconfined concrete strain at maximum strength}
}
\newglossaryentry{esp}
{
    type=matprop,
    name=$\epsilon_{sp}$,
    description={unconfined concrete crushing (spalling) strain}
}
\newglossaryentry{Ec}
{
    type=matprop,
    name=$E_{c}$,
    description={concrete compressive elastic modulus}
}
\newglossaryentry{nuc}
{
    type=matprop,
    name=$\nu_{c}$,
    description={concrete Poisson's ratio}
}
\newglossaryentry{Gc}
{
    type=matprop,
    name=$G_{c}$,
    description={concrete shear modulus}
}
\newglossaryentry{Es}
{
    type=matprop,
    name=$E_{s}$,
    description={steel tensile elastic modulus}
}
\newglossaryentry{Esh}
{
    type=matprop,
    name=$E_{sh}$,
    description={steel initial strain hardening tangent}
}
\newglossaryentry{fy}
{
    type=matprop,
    name=$f_{y}$,
    description={steel yield strength}
}
\newglossaryentry{fu}
{
    type=matprop,
    name=$f_{u}$,
    description={steel ultimate strength}
}
\newglossaryentry{esh}
{
    type=matprop,
    name=$\epsilon_{sh}$,
    description={steel strain at initial strain hardening}
}
\newglossaryentry{esu}
{
    type=matprop,
    name=$\epsilon_{su}$,
    description={steel strain at peak stress}
}
\newglossaryentry{et}
{
    type=matprop,
    name=$\epsilon_{t}$,
    description={unconfined concrete tensile cracking strain}
}

\newglossaryentry{f2e}
{
    type=modelprop,
    name=$f_{2e}$,
    description={concrete effective confinement strength}
}
\newglossaryentry{fcc}
{
    type=modelprop,
    name=$f'_{cc}$,
    description={concrete peak compressive strength}
}
\newglossaryentry{ecc}
{
    type=modelprop,
    name=$\epsilon_{cc}$,
    description={concrete strain at peak strength}
}
\newglossaryentry{ecu0}
{
    type=modelprop,
    name=$\epsilon_{cu0}$,
    description={Mander concrete crushing strain}
}
\newglossaryentry{xi}
{
    type=modelprop,
    name=$\xi$,
    description={Mander concrete parameter}
}
\newglossaryentry{r}
{
    type=modelprop,
    name=$r$,
    description={Mander concrete parameter}
}
\newglossaryentry{fcu0}
{
    type=modelprop,
    name=$f_{cu0}$,
    description={Mander concrete crushing strength}
}
\newglossaryentry{fcu}
{
    type=modelprop,
    name=$f_{cu}$,
    description={concrete crushing strength}
}
\newglossaryentry{ecu}
{
    type=modelprop,
    name=$\epsilon_{cu}$,
    description={concrete crushing strain}
}
% \begin{figure}
%     \centering
%     \scalebox{0.75}{\input{figures/plasticity.tikz}}
%     \caption{The structure of plasticity theories.}
%     \label{fig:plasticity}
% \end{figure}


% Five key concepts form the basis of almost all classical theories of
% plasticity. They are

% \begin{enumerate}
% \def\labelenumi{\arabic{enumi}.}
% \tightlist
% \item
%   The decomposition of strain into elastic and plastic parts;
% \item
%   Yield criteria, which predict whether the solid responds elastically
%   or plastically;
% \item
%   Strain hardening rules, which control the way in which resistance to
%   plastic flow increases with plastic straining;
% \item
%   The plastic flow rule, which determines the relationship between
%   stress and plastic strain under multi-axial loading;
% \item
%   The elastic unloading criterion, which models the irreversible
%   behavior of the solid.
% \end{enumerate}

\[
\dot{\bar{\epsilon}}^p \triangleq \left|\dot{\epsilon}^p\right|
\qquad\text{ so }\qquad
\bar{\epsilon}^p(t) = \int_0^t \dot{\bar{\epsilon}}^p(\zeta) d \zeta
\]
\begin{marsdenbox}
    \begin{enumerate}
\def\labelenumi{\arabic{enumi}.}
\item Kinematic decomposition
\[
\boldsymbol{\varepsilon} = \boldsymbol{\varepsilon}^e + \boldsymbol{\varepsilon}^p
\]
\item
  Elastic stress-strain relationships \[
    \begin{aligned}
    \boldsymbol{\sigma} &=\frac{\partial W\left(\varepsilon-\varepsilon^{p}\right)}{\partial \varepsilon}=\mathbf{C}\left[\varepsilon-\varepsilon^{p}\right] \\
    %\mathbf{C}: &=\frac{\partial^{2} W\left(\varepsilon-\varepsilon^{p}\right)}{\partial \varepsilon^{2}}=\text { constant (elastic moduli). }
    \end{aligned}
    \]
\item
  Elastic domain in stress space (single surface) \[
  \mathbb{E}_{\sigma}=\left\{(\boldsymbol{\sigma}, \boldsymbol{q}) \in \mathbb{S} \times \mathbb{R}^{m} \mid f(\boldsymbol{\sigma}, \boldsymbol{q}) \leq 0\right\}
  \]
  where $f$ is such that:
  $$
  \text { if } f=0 \text {, then } \dot{f} \leq 0
  $$
\item
  Flow rule and hardening law

  \[
  \begin{aligned}
  & \text{General}  & \text{Associative} \\
  \dot{\varepsilon}^{p} &=\lambda \boldsymbol{r}(\boldsymbol{\sigma}, \boldsymbol{q}) &=\lambda \frac{\partial f}{\partial \boldsymbol{\sigma}}  \\
  \dot{\boldsymbol{q}} &=-\lambda \boldsymbol{h}(\boldsymbol{\sigma}, \boldsymbol{q}) &= -\lambda \mathbf{D} \frac{\partial f}{\partial \boldsymbol{q}}
  \end{aligned}
  \]
  where \(\mathbf{D}\) is a matrix of generalized plastic moduli.
\item
  Kuhn-Tucker complementarity (loading/unloading)
  
\[
  \lambda \geq 0, \quad f(\boldsymbol{\sigma}, \boldsymbol{q}) \leq 0, \quad \lambda f(\boldsymbol{\sigma}, \boldsymbol{q})=0 .
\]

\item
  Consistency condition: plastic flow ($\dot{\bar{\epsilon}}^p>0$) can only occur if the state of stress persists in the yield set (ie,  remains an elastic-plastic state)
    $$
    \text { if } \dot{\epsilon}^{\mathrm{p}} \neq 0 \text {, then } f=0 \text { and } \dot{f}=0
    $$ 
% or
%   \[
%   \lambda \dot{f}(\boldsymbol{\sigma}, \boldsymbol{q})=0 .
%   \]
\end{enumerate}
\end{marsdenbox}


\iffalse
\[
H\left(\bar{\epsilon}^p\right) \triangleq \frac{d Y\left(\bar{\epsilon}^p\right)}{d \bar{\epsilon}^p}
\]

\[\dot{\epsilon}^p=\left\{\begin{array}{llll}0 & \text { if } & |\sigma|<Y\left(\bar{\epsilon}^p\right) & \text { (behavior in the elastic state), } \\ 0 & \text { if } & |\sigma|=Y\left(\bar{\epsilon}^p\right) & \text { and } \quad \sigma \dot{\epsilon}<0 \quad \text { (elastic unloading), } \\ \beta \dot{\epsilon} & \text { if } & |\sigma|=Y\left(\bar{\epsilon}^p\right) & \text { and } \quad \sigma \dot{\epsilon}>0 \quad \text { (plastic loading), }\end{array}\right.
\]
\fi

\begin{longtable}[]{@{}lll@{}}
\toprule\noalign{}
& Govindjee & Simo \\
\midrule\noalign{}
\endhead
\bottomrule\noalign{}
\endlastfoot
\(\sigma_\text{yield}\) & \(Y\) & \(K\) \\
\(H_\text{iso}\) & \(H_\text{iso}\) & \(K^\prime\) \\
\(H_\text{kin}\) & & \(H^\prime\) \\
\(\dot{\sigma}_y\) & \(\dot{Y}\) & \(K^\prime\dot{\alpha}\) \\
\(\|\epsilon^p\|\) & \(\bar{\epsilon}^p\) & \(\alpha\) \\
& \(\sigma_\text{eff}\) & \(\eta\) \\
\(\sigma_\text{back}\) & & \\
\(\sigma_\text{back}\) & & \\
\end{longtable}

% \hypertarget{yield-criteria}{%
% \paragraph{Yield Criteria}\label{yield-criteria}}

% \begin{description}
% \tightlist
% \item[Huber Von-Mises (\(J_2\))]
% \end{description}

% \begin{quote}
% \[
% \begin{aligned}
% \boldsymbol{\eta} &\triangleq \operatorname{dev}[\boldsymbol{\sigma}]-\overline{\boldsymbol{\beta}} \\
% \operatorname{tr}[\overline{\boldsymbol{\beta}}] &\triangleq 0 \\
% f(\boldsymbol{\sigma}, \boldsymbol{q}) &=\|\boldsymbol{\eta}\|-\sqrt{\frac{2}{3}} K(\alpha) \\
% \dot{\varepsilon}^{p} &=\lambda \frac{\boldsymbol{\eta}}{\|\boldsymbol{\eta}\|}, \\
% \lambda &= \left\|\dot{\varepsilon}^{p}\right\|=\frac{\left<\frac{\boldsymbol{\eta}}{\|\boldsymbol{\eta}\|}: \dot{\boldsymbol{\varepsilon}}\right>}{1+\frac{H^{\prime}+K^{\prime}}{3 \mu}}, \\
% \dot{\overline{\boldsymbol{\beta}}} &=\lambda \frac{2}{3} H^{\prime}(\alpha) \frac{\boldsymbol{\eta}}{\|\boldsymbol{\eta}\|} \\
% \dot{\alpha} &=\lambda \sqrt{\frac{2}{3}} \\
% % \boldsymbol{n}&:=\frac{\boldsymbol{\eta}}{\|\boldsymbol{\eta}\|}
% \end{aligned}
% \]
% \end{quote}

\hypertarget{isotropic-hardening}{%
\subsection{Isotropic Hardening}\label{isotropic-hardening}}
\begin{description}
\item[Prager]



\item[Voce] (Voce 1948). 
  \[
  \boxed{
  \dot{\sigma}_\text{y} = \dot{\lambda} H_\text{iso}\left(1-\frac{\sigma_\text{y}}{f_{\text{y}, \infty}}\right)
  }
  \] 

% \begin{quote}
\[
\sigma_{y} = f_{y} + f_{\infty}\left(1-e^{-b \varepsilon_{e q}^{p}}\right)
%Y\left(\bar{\varepsilon}^{\mathrm{p}}\right)=f_{s}-\left(Y_{s}-Y_{0}\right) \exp \left(-\frac{H_{0}}{Y_{s}} \bar{\varepsilon}^{\mathrm{p}}\right)
\]

\begin{longtable}[]{@{}ll@{}}
\toprule\noalign{}
Symbol & Description \\
\midrule\noalign{}
\endhead
\bottomrule\noalign{}
\endlastfoot
\(f_y\) & Initial yield stress. \\
\(f_\infty\) & Saturated yield stress \\
\end{longtable}
% \end{quote}


% \hypertarget{power-law-hardening-govindjee-2020}{%
% \subparagraph{Power-Law Hardening (Govindjee
% 2020)}\label{power-law-hardening-govindjee-2020}}
\item[Power-Law Hardening] The aforementioned hardening law can be considered a specialization of the more general rule:
\[
\boxed{
\dot{\sigma}_y(\bar{{\varepsilon}}^\mathrm{p})=\dot{\bar{\epsilon}}^p H_\text{iso}\left(1-\frac{\sigma_y}{f_{y,s}}\right)^{r}, \\
\quad \sigma_y(0) = Y_0
}
\]
Note that this IVP can be integrated in closed form, furnishing the solution: \[
\sigma_y\left(\bar{{\varepsilon}}^{\mathrm{p}}\right)
= \begin{cases}
f_{s}-\left[\left(f_{y,s}-f_{y}\right)^{(1-r)}+(r-1)\left(\frac{H_{0}}{\left(f_{s}\right)^{r}}\right) \bar{\varepsilon}^{\mathrm{p}}\right]^{\frac{1}{1-r}} & r \ne 1 \\
f_{s}-\left(f_{y,s}-f_y\right) \exp \left(-\frac{H_\text{iso}}{f_{s}} \bar{{\varepsilon}}^{\mathrm{p}}\right) & r = 1
\end{cases}
\]
\end{description}

\hypertarget{kinematic-hardening}{%
\subsection{Kinematic Hardening}\label{kinematic-hardening}}

\begin{description}
\item[Prager] Melan (1938), Ishlinskii (1954), and Prager (1956).
  % \begin{framed}
    \[
    \sigma_{\mathrm{b}} = H_\mathrm{k} \, \varepsilon^p
    \]
    \begin{longtable}[]{@{}ll@{}}
    \toprule\noalign{}
    Symbol & Description \\
    \midrule\noalign{}
    \endhead
    \bottomrule\noalign{}
    \endlastfoot
    \({H}_{\text{k}}\) & Kinematic hardening modulus \\
    \end{longtable}
  % \end{framed}
\item[Ziegler]
    \[
    d \alpha_{i j}=d a\left(\varepsilon_{i j}^p\right)\left(\sigma_{i j}-\alpha_{i j}\right)
    \]
    
    for some scalar function \(a\). or
    
    \[
    \dot{\boldsymbol{\alpha}}=C \dot{\bar{\varepsilon}}^{p l} \frac{1}{\sigma^0}(\boldsymbol{\sigma}-\boldsymbol{\alpha})+\frac{1}{C} \boldsymbol{\alpha} \dot{C}
    \] 
    This form of evolution law for
    \(\alpha\) defines the rate of \(\alpha\) due to plastic straining to be
    in the direction of the current radius vector from the center of the
    yield surface, \(\sigma-\alpha\), and the rate due to temperature
    changes to be toward the origin of stress space. Rice
    (1975) writes this concept quite generally as 
    \[
    \dot{\boldsymbol{\alpha}}=\dot{\mu}(\boldsymbol{\sigma}-\boldsymbol{\alpha})+h \boldsymbol{\alpha} \dot{\theta} .
    \]
    
    % \begin{quote}
    \[
    \dot{\boldsymbol{\sigma}}_\text{back}
    =\frac{2}{3} H \dot{\varepsilon}^{p}
    =\dot{\lambda} \sqrt{\frac{2}{3}} H \frac{\boldsymbol{\eta}}{\|\eta\|}
    \]
    \begin{longtable}[]{@{}ll@{}}
    \toprule\noalign{}
    Symbol & Description \\
    \midrule\noalign{}
    \endhead
    \bottomrule\noalign{}
    \endlastfoot
    \(\bar{H}_{\text{kin}}\) & Kinematic hardening modulus \\
    \end{longtable}
    % \end{quote}

\item[Armstrong--Frederick]

\begin{quote}
\[
\begin{aligned}
\dot{\boldsymbol{\sigma}}_\text{back} &=\frac{2}{3} H_\text{k} \dot{\boldsymbol{\varepsilon}}^{\text{p}}-\dot{\bar{\varepsilon}}^{\text{p}} \gamma \boldsymbol{\sigma_\text{back}} \\
%&=\dot{\lambda}\left(\frac{2}{3} \bar{H}_\text{k} \frac{\partial \Phi}{\partial \boldsymbol{\sigma}}-b \boldsymbol{\sigma_\text{back}}\right)
\end{aligned}
\]

\begin{longtable}[]{@{}ll@{}}
\toprule\noalign{}
Symbol & Description \\
\midrule\noalign{}
\endhead
\bottomrule\noalign{}
\endlastfoot
\(H_{\text{k}}\) & Back stress modulus \\
\(\gamma\) & Kinematic saturation parameter \\
\end{longtable}
\end{quote}

% \item[Chaboche]

% % \hypertarget{chaboche-chaboche-et-al.-1979}{%
% % \subparagraph{Chaboche (Chaboche et al.~1979)}\label{chaboche-chaboche-et-al.-1979}}

% \[
% \boxed{
% \begin{gathered}
% \dot{\boldsymbol{\sigma}}_{\text{back},k}=\sqrt{2 / 3} C_{k} \dot{\boldsymbol{\varepsilon}}^{p} -\gamma_{k} \dot{\bar{\varepsilon}}^{p} \boldsymbol{\sigma}_{\text{back},k}\\
% \mathbf{\sigma}_\text{back}=\sum_{k=1}^{n B a c k} \sigma_{\text{back},k}
% \end{gathered}
% }
% \]
% \[
% \dot{\boldsymbol{\alpha}}=C \dot{\bar{\varepsilon}}^{p l} \frac{1}{\sigma^0}(\boldsymbol{\sigma}-\boldsymbol{\alpha})-\gamma \boldsymbol{\alpha} \dot{\bar{\varepsilon}}^{p l}+\frac{1}{C} \boldsymbol{\alpha} \dot{C}
% \] where \(C\) and \(\gamma\) are material parameters,

\end{description}

Ziegler hardening is typically implemented with isotropic hardening as:

\[
\begin{aligned}
H^{\prime}(\alpha) &=(1-\theta) \bar{H} \\
K(\alpha) &=\left[\sigma_{y}+\theta \bar{H} \alpha\right], \quad \theta \in[0,1]
\end{aligned}
\]
